\section{מבוא לסוכני AI}

בעשור האחרון עברה הבינה המלאכותית ממודלים של סיווג וחיזוי סטטיסטי, אל מערכות
שמסוגלות לנהל שיחה, לתכנן פעולות ולבצע משימות מורכבות בסביבה דיגיטלית.
במסגרת קורס \eng{AI Agents} נלמד כיצד לבנות, להפעיל ולנהל סוכנים --- ישויות
תוכנה שמשלבות מודל שפה גדול (\eng{LLM}) עם זיכרון עבודה, כלים חיצוניים ויכולת
תכנון. מסמך זה מסכם את הארכיטקטורה כפי שנלמדה בהרצאה 07, ומרחיב אותה לכיוון
יישומי בסייבר.

\subsection{מהו סוכן AI?}

סוכן \eng{AI} הוא מערכת שמקבלת מטרה בשפה טבעית, מפרקת אותה לצעדים, בוחרת כלים
מתאימים ומחזירה תוצאה. בניגוד לצ'אטבוט קלאסי, הסוכן \emph{פועל}: הוא יכול
לקרוא קבצים, להפעיל \eng{API}, לחפש במאגר ידע או לקמפל מסמך \LaTeX{} ל-\eng{PDF}.
לפי חומר ההרצאה, הסוכן אינו עומד בפני עצמו --- הוא נשען על שכבות תחתונות של
תוכנה קבועה ושכבת \eng{Skill} שמעטפת את הקוד בשפה טבעית \cite{brown2020gpt3}.

\begin{insightbox}
המעבר ממודל שפה לסוכן פעיל דורש שלושה מרכיבים: \textbf{תכנון} (מה לעשות),
\textbf{ביצוע} (כלים) ו\textbf{זיכרון} (הקשר + אחסון חיצוני). ללא אחד מהם,
המערכת נשארת מחולל טקסט בלבד.
\end{insightbox}

\subsection{מדוע ארכיטקטורה חשובה?}

ארגונים שמטמיעים סוכני \eng{AI} נתקלים במהירות בבעיות חוזרות: עלויות טוקנים
גבוהות, תוצאות לא עקביות, קושי בבקרת איכות ותלות ב-\eng{GUI} במקום ב-\eng{API}.
הבנת הארכיטקטורה --- שלוש השכבות, \eng{RAG}, \eng{MCP}, תבניות \eng{CLS}
ורב-סוכנים --- מאפשרת לבנות מערכות \textbf{רפטביליות} (\eng{reproducible})
שחוזרות על אותה איכות בכל הרצה.

\subsection{מבנה המסמך}

המסמך בנוי מ-15 פרקים: מיסודות הסוכן ושכבותיו, דרך מנגנוני ידע ותקשורת,
ועד יצירת \eng{PDF} מקצועי, בקרת איכות ואורכוסטרציה. פרק 13 מדגים יישום
קצר בתחום \eng{CTI}/\eng{OSINT} וחקירת \eng{Ransomware}. הפרקים האחרונים
מסכמים את המשמעות הניהולית והטכנולוגית.

\subsection{מטרות למידה}

לאחר קריאת המסמך, הקורא אמור להבין:
\begin{itemize}[rightmargin=2em]
  \item את ההבדל בין שכבת תוכנה, \eng{Skill} וסוכן.
  \item כיצד \eng{RAG} מרחיב את חלון ההקשר מבלי לאמן מחדש את המודל.
  \item מדוע \eng{MCP} שונה מ-\eng{Skill} סטטי עם \eng{API}.
  \item מדוע עבודה טקסטואלית/\eng{API} עדיפה על סוכן דפדפן לייצור מסמכים.
  \item כיצד תבנית \eng{CLS} וסקילי \eng{QA} מבטיחים רפטביליות.
\end{itemize}

\begin{notebox}
מונחים באנגלית מסומנים בפקודה \texttt{\textbackslash eng\{\}} --- לדוגמה:
\eng{Context Window}, \eng{Retrieval-Augmented Generation}.
\end{notebox}

\subsection{הקשר אקדמי}

מטלה 04 בקורס \eng{AI Agents} דורשת הפקת מסמך \eng{PDF} מקצועי באמצעות
\eng{LuaLaTeX} וקובץ \eng{CLS} מותאם. המסמך משקף את יכולת הסטודנט להגדיר
הנחיות לסוכן, לבנות תבנית עיצוב נפרדת ולייצר תוכן אקדמי מלא --- לא שלד ריק.
המקור העיקרי לתוכן הוא תמלול הרצאה 07 (\texttt{sources/transcript-lecture07.txt}).
